\section{A coherent seed at quadratic logarithmic-domain cost}
\label{sec:canonical-spike-prefix}
Write $F(q)=\E\tanh^2(\sqrt qG)$ and $\chi(q)=1-F(q)$. A covariance spike of size $1+\alpha$ evolves variance and directional correlation by
\begin{equation}
 v=q+\alpha r^2,\qquad q'=F(v),\qquad
 r'=\sqrt{1+\alpha}\,r\chi(v).
 \label{eq:canonical-spike-update}
\end{equation}
For a valid state $0<r^2\le q\le1$, its coherence is $\theta=r^2/q$.

\begin{lemma}[Native critical decay preserves coherence up to a constant]
\label{lem:native-coherence}
There is an absolute $C$ such that any number of zero-spike transitions satisfy $\theta_j\ge e^{-C}\theta_1$.
\end{lemma}
\begin{proof}
For a zero spike the coherence ratio is
\[
 L(q)=\frac{q\chi(q)^2}{F(q)}.
\]
Stein and Cauchy--Schwarz give $0<L(q)\le1$. Taylor expansion at zero gives
\[
 F(q)=q-2q^2+\tfrac{17}{3}q^3+O(q^4),\qquad
 L(q)=1-\tfrac23q^2+O(q^3).
\]
Thus $-\log L(q)/(q-F(q))$ has finite limit $1/3$ at zero. It is continuous and finite on $(0,1]$, so $-\log L(q)\le C(q-F(q))$ there. The variance follows $q'=F(q)$; summing this inequality telescopes to at most $Cq_1\le C$.
\end{proof}

\begin{lemma}[A constant-coherence prefix]
\label{lem:spike-prefix}
Fix $C_0,c_0>0$. There are absolute positive constants $K,C,c$ depending only on them such that the following holds for all sufficiently large $H$ and $1\le D\le\sqrt H$. Suppose
\[
 c_0/(D+1)\le q_1\le1,\qquad
 e^{-C_0(D+1)}\le|r_1|\le\sqrt{q_1}/2.
\]
Set $\alpha=K(D+1)/H$. Use zero spikes for the first $\lfloor H/4\rfloor$ transitions and then constant spike $\alpha$. At every step between $\lceil3H/4\rceil$ and $H$,
\[
 c\alpha\le q_\ell\le C\alpha,
 \qquad r_\ell^2/q_\ell\ge1/2.
\]
The covariance entropy action is at most $K^2(D+1)^2/(2H)$. The fixed constant $K$ may be enlarged as needed, after increasing the depth threshold.
\end{lemma}
\begin{proof}
Native reciprocal variance estimates give $c/H\le q\le C/H$ at the end of the initial burn, using $q_1\ge c_0/(D+1)$ and $D\le\sqrt H$. Lemma~\ref{lem:native-coherence} gives coherence at least $e^{-C_1(D+1)}$ there; no logarithmic loss in $H$ occurs.

For the subsequent constant spike, $r^2\le q$ and the quadratic drift imply
\[
 q'\le F((1+\alpha)q)\le(1+\alpha)q-cq^2.
\]
Starting from $q\le C/H$, this preserves $q\le C_2\alpha$ after choosing a large enough fixed $C_2$. The native lower comparison $q'\ge F(q)$ also ensures $q\ge c/H$ over the full horizon.

The exact coherence update is
\[
 \theta'=\frac{(1+\alpha)\theta}{1+\alpha\theta}
             L(q(1+\alpha\theta)).
\]
Since $L(v)\ge1-Cv^2$ at small $v$ and $q\le C_2\alpha$, while $\theta\le1/2$ this gives $\theta'\ge(1+\alpha/4)\theta$ for small enough $\alpha$. It reaches $1/2$ within $C_3(D+1)/\alpha=C_3H/K$ steps. Thereafter the boundary $1/2$ is invariant, because the logistic increment there is of order $\alpha$ and dominates the order-$\alpha^2$ loss.

Once $\theta\ge1/2$, whenever $q\le c_1\alpha$ for sufficiently small $c_1>0$,
\[
 q'=F(q(1+\alpha\theta))\ge q(1+\alpha/4).
\]
It reaches $c_1\alpha$ from $c/H$ within
$C_4\alpha^{-1}\log(C_5\alpha H)$ steps. This is at most
$C_6H\log(C_7K)/K$, uniformly in $D\ge1$, because
$\log(C_5K(D+1))/(D+1)\le C\log(C_7K)$.
The lower boundary is invariant by monotonicity of $F$. Choose the fixed $K$ so that the two growth periods total less than $H/4$. The claimed state bounds then hold by step $3H/4$ and remain valid. Finally $\alpha-\log(1+\alpha)\le\alpha^2/2$ proves the action bound.
\end{proof}

\begin{lemma}[Exact orthogonal-defect identity]
\label{lem:spike-defect}
For any nonnegative spike, put $e=q-r^2$. Then
\[
 e'=J(q')+e(1-q')^2,
 \qquad J(u)=u-F^{-1}(u)(1-u)^2.
\]
For small $u\ge0$, $J(u)=\tfrac23u^3+O(u^4)$ and $0\le J(u)\le Cu^3$.
\end{lemma}
\begin{proof}
Since $(1+\alpha)r^2=v-e$ and $q'=F(v)$, the displayed recurrence gives $r'^2=(v-e)(1-q')^2$. Taylor expansion yields $F^{-1}(u)=u+2u^2+(7/3)u^3+O(u^4)$ and hence the expansion for $J$. Nonnegativity follows from $v(\E\tanh'(\sqrt vG))^2\le\E\tanh^2(\sqrt vG)$.
\end{proof}
